% QPAD — CS 228 Checkpoint 1 report (IEEE conference, double column).
% Compile on Overleaf (pdfLaTeX + BibTeX). Every number comes from tables/generated_values.tex, which is
% produced by scripts/build_report_assets.py from result files. Unmeasured values appear as
% [INSERT MEASURED RESULT FROM ...] — never replace them by hand with invented values.
% Scope changes since the proposal are marked with \scope{...} (blue).
\documentclass[conference]{IEEEtran}
\usepackage[utf8]{inputenc}
\usepackage[T1]{fontenc}
\usepackage{graphicx}
\usepackage{booktabs}
\usepackage{amsmath}
\usepackage{xcolor}
\usepackage{url}
\usepackage[hidelinks]{hyperref}
\usepackage{balance}   % equal column heights on the last page

\newcommand{\scope}[1]{\textcolor{blue}{#1}}
% Many floats in a short two-column paper: allow more per page so they stay near their text.
\renewcommand{\topfraction}{0.9}\renewcommand{\bottomfraction}{0.6}\renewcommand{\textfraction}{0.08}
\renewcommand{\floatpagefraction}{0.8}\setcounter{topnumber}{3}\setcounter{bottomnumber}{2}
\setcounter{totalnumber}{5}
\newcommand{\todo}[1]{\textbf{\textcolor{red}{[#1]}}}
% Figure that compiles before the asset exists.
\newcommand{\qpadfig}[3]{%
  \begin{figure}[!tb]\centering
  \IfFileExists{figures/#1.pdf}{\includegraphics[width=\columnwidth]{figures/#1.pdf}}%
  {\fbox{\parbox{0.9\columnwidth}{\centering\vspace{1.2cm}\todo{INSERT FIGURE figures/#1.pdf FROM results/}\vspace{1.2cm}}}}
  \caption{#2}\label{#3}\end{figure}}
\newcommand{\qpadtable}[1]{\IfFileExists{tables/#1.tex}{\input{tables/#1.tex}}{\todo{INSERT tables/#1.tex}}}

% AUTO-GENERATED by qpad.reporting.build_report_assets — do not edit by hand.
\newcommand{\QPADAuditStatus}{PASS}
\newcommand{\QPADBestDevLoss}{0.0532}
\newcommand{\QPADBestEpoch}{20}
\newcommand{\QPADDevACER}{0.37\%}
\newcommand{\QPADDevACERmax}{0.56\%}
\newcommand{\QPADDevAPCER}{2/540 = 0.37\%}
\newcommand{\QPADDevAPCERatThr}{2/540}
\newcommand{\QPADDevAPCERmax}{2/270 = 0.74\%}
\newcommand{\QPADDevAUROC}{0.9998}
\newcommand{\QPADDevBPCER}{1/270 = 0.37\%}
\newcommand{\QPADDevBPCERatThr}{1/270}
\newcommand{\QPADDevEER}{0.37\%}
\newcommand{\QPADDevPrintAPCER}{2/270 = 0.74\%}
\newcommand{\QPADDevReplayAPCER}{0/270 = 0.00\%}
\newcommand{\QPADDuration}{3.10--5.06\,s (median 4.97\,s)}
\newcommand{\QPADEpochsRun}{25}
\newcommand{\QPADFps}{29.84}
\newcommand{\QPADFrameCount}{41--151 (median 146)}
\newcommand{\QPADGitCommit}{b9a416a5dd96}
\newcommand{\QPADNumFailed}{3 (6\_1\_12\_1, 3\_1\_54\_5, 5\_1\_54\_3)}
\newcommand{\QPADNumVideoMembers}{4950}
\newcommand{\QPADNumVideos}{2970}
\newcommand{\QPADResolutions}{1080x1920 (2970)}
\newcommand{\QPADTestACER}{3.12\%}
\newcommand{\QPADTestACERmax}{4.72\%}
\newcommand{\QPADTestAPCER}{27/720 = 3.75\%}
\newcommand{\QPADTestAPCERmax}{25/360 = 6.94\%}
\newcommand{\QPADTestAUROC}{0.9963}
\newcommand{\QPADTestBPCER}{9/360 = 2.50\%}
\newcommand{\QPADTestEER}{2.78\%}
\newcommand{\QPADTestPrintAPCER}{2/360 = 0.56\%}
\newcommand{\QPADTestReplayAPCER}{25/360 = 6.94\%}
\newcommand{\QPADThreshold}{0.731701}


\begin{document}

\title{QPAD: Calibration-Set Composition Bias in INT8\\ Face Presentation Attack Detection}

\author{\IEEEauthorblockN{Ananya Tyagi \quad Vindhya}
\IEEEauthorblockA{CS 228 Biometric Security with AI, San Jos\'e State University\\
\textit{Nature of project:} Test whether INT8 calibration-set composition creates attack-specific regressions
in face presentation attack detection.}}

\maketitle

\begin{abstract}
\scope{We study whether the class mixture of the post-training INT8 calibration set changes how often a face
presentation attack detector (DeepPixBiS) accepts print or replay attacks on OULU-NPU Protocol~II.}
This checkpoint reports the data engineering pipeline, the dataset audit, exploratory analysis and the FP32
baseline with a development-set operating threshold. Results are reported at commit \texttt{\QPADGitCommit}.
\end{abstract}

% =====================================================================================================
\section{Problem Statement and Motivation}
A face presentation attack detection (PAD) system decides whether a camera sees a live person
(\emph{bona fide}) or a presentation attack instrument (PAI) such as a printed photograph or a video replayed on a
screen. A missed attack is a security failure: an impostor holding a photo or a phone is accepted as the
enrolled user. Errors are therefore reported separately as APCER (attacks accepted) and BPCER (bona fide
presentations rejected) following ISO/IEC 30107-3~\cite{iso30107-3}.

Deployed PAD models are often quantized to INT8 for speed~\cite{jacob2018quantization,migacz2017int8}.
Post-training INT8 quantization keeps the trained weights but selects per-tensor value ranges from a small
\emph{calibration set}. \scope{If that set over- or under-represents an attack type, the chosen ranges might clip
exactly the activations that distinguish that attack, raising its APCER even though no weight changes.}

\textbf{Research question.} \scope{Does the composition of the INT8 calibration set cause attack-specific
degradation---in print or replay APCER---relative to an INT8 engine calibrated on the natural class mixture,
after every engine receives its own development-set threshold?} The primary comparison is \emph{unusual-calibration
INT8 versus natural-calibration INT8 of the same nested seed}, which isolates calibration composition from
ordinary quantization error. A meaningful effect is pre-registered as an attack-specific APCER increase of
5 percentage points; a null result with tight confidence intervals is an acceptable, informative outcome.

% =====================================================================================================
\section{Data Engineering}
\subsection{Dataset origin, access and characteristics}
\scope{We use OULU-NPU~\cite{boulkenafet2017oulu}}, a mobile face PAD database of 55 subjects recorded with the
front cameras of six smartphones in three sessions with different illumination and background, with print and
video-replay attacks produced by two printers and two displays~\cite{boulkenafet2017oulu}. OULU-NPU is distributed
for research use under an end-user licence agreement; our copy was provided by the course instructor on request
for this course project. The data are kept in private storage and are not redistributed, and no face images appear
in this report.

Our archive audit found \QPADNumVideoMembers{} video files in the three split archives, of which
\QPADNumVideos{} belong to Protocol~II. Protocol~II videos have resolution \QPADResolutions{}
(width\,$\times$\,height, number of videos), \QPADFrameCount{} frames, a median of \QPADFps{} frames/s and
duration \QPADDuration{}.
Each video has a same-name text file with per-frame eye coordinates (\texttt{frame, x1, y1, x2, y2}).

\subsection{Protocol~II splits and labels}
Filenames encode \texttt{phone\_session\_subject\_PAI}. PAI code 1 is bona fide, 2--3 print and 4--5 replay.
Protocol~II trains and selects thresholds on Printer~1/Display~1 attacks (codes 2, 4) and tests on unseen
Printer~2/Display~2 attacks (codes 3, 5); subjects 1--20, 21--35 and 36--55 form disjoint train, development
and test splits (Table~\ref{tab:counts}). The development split is used only for model selection and threshold
choice. The test split is scored only after the threshold is frozen; a lock file records the threshold's
hash and refuses any later test evaluation with a different threshold.

\begin{table}[!tb]\centering\caption{Protocol II videos per split and class (from the audit manifest).}
\label{tab:counts}\qpadtable{class_counts}\end{table}

\subsection{Schema, quality checks and cleaning}
One manifest row per video records split, subject, phone, session, PAI code, class, frame count, resolution,
FPS, duration, annotation statistics, decode status and file size. Overall audit status:
\textbf{\QPADAuditStatus}. The audit verifies: the 2,970 expected videos exist with the expected split/class
counts; the subject/PAI membership rule agrees with the official Protocol~II list files; no subject appears in
two splits; there are no duplicate or unparseable identifiers; every video decodes completely with FFmpeg; every
annotation file is well formed; and every video, apart from the two documented test exceptions below, has at
least 20 valid annotated frames. Annotation rows whose
coordinates are all zero (missing eye detections) are treated as invalid frames, not imputed.
Videos failing at least one per-video check: \QPADNumFailed{}. All 2,970 videos are kept, with a documented
rule fixed before any training or test scoring: one training video has a single damaged MJPEG frame, which is
excluded from frame selection; two test attack videos have only 2 and 19 frames with valid eye coordinates and
are scored on all of those frames instead of 20. We keep rather than drop them because removing attacks that are
hard to annotate could make APCER look better than it is.

\subsection{Preprocessing}
For each video we select 20 evenly spaced frames among frames with valid eye annotations and save their indices.
Using the supplied eye coordinates (no separate detector), an exact similarity transform places the two eyes at
fixed positions of a $224\times224$ crop; frames are converted to RGB and normalized with ImageNet statistics
at load time, identically for every backend. \scope{For calibration, five evenly spaced frames of each training
video are saved.} No feature-reduction method such as PCA is used, because DeepPixBiS learns image features
directly from pixels. Processed crops are stored in verified, compressed shards in private storage.

\subsection{Shortcomings and limitations}
Protocol~II covers two printers and two displays only; subjects are few (20 test subjects), so test videos are
clustered by subject and we resample subjects, not videos, for uncertainty. \scope{Because test attacks use
equipment unseen in training (Printer~2, Display~2), a ``print'' or ``replay'' effect is also an effect on that
specific instrument. Each calibration set is 120 videos $\times$ 5 frames; frames of one video are correlated, so
the effective calibration sample is closer to 120 than to 600.} Eye annotations are used as supplied. Conclusions are limited to OULU-NPU Protocol~II.

% =====================================================================================================
\section{Exploratory Data Analysis}
Fig.~\ref{fig:balance} shows class balance: attacks outnumber bona fide videos two to one, so a single accuracy
would hide attack-specific errors; we always report per-class rates. Fig.~\ref{fig:frames} is a univariate view
of crop sharpness, and Fig.~\ref{fig:scatter} a bivariate view of crop brightness versus sharpness per video.
Pixel statistics use train and development videos only.
\qpadfig{class_balance_by_split}{Protocol~II videos per split and class (class balance).}{fig:balance}
\qpadfig{sharpness_by_class}{Univariate: per-video sharpness (variance of Laplacian) of aligned crops by class
(train + development).}{fig:frames}
\qpadfig{brightness_vs_sharpness}{Bivariate: brightness versus sharpness, one point per video, coloured by
class (train + development).}{fig:scatter}
Within every split the three classes are exactly balanced (Table~\ref{tab:counts}), and every
phone and session contributes the same number of videos per class, so no device or recording condition dominates.
On the aligned crops of the 1,890 train and development videos (630 per class), attacks are brighter than bona fide
presentations (median per-video grey level 109.7 for bona fide, 130.8 for print and 144.7 for replay), and print
attacks are less sharp (median variance of Laplacian 2186, versus 3568 for bona fide and 2955 for replay;
Fig.~\ref{fig:frames}), consistent with the loss of high-frequency detail when a face is printed and re-captured.
The three classes nevertheless overlap heavily in the brightness--sharpness plane (Fig.~\ref{fig:scatter}), so
these global image statistics cannot separate attacks on their own; the detector has to learn finer, local cues,
which motivates a pixel-wise model.

% =====================================================================================================
\section{Baseline Implementation and Preliminary Results}
\subsection{Model and training}
\scope{DeepPixBiS~\cite{george2019deep}} uses the first eight blocks of an ImageNet-pretrained
DenseNet-161~\cite{huang2017densely}, a $1\times1$ convolution with sigmoid giving a $14\times14$ realness map, and
a linear layer giving a binary output; the loss is $0.5$ pixel-wise BCE $+\,0.5$ binary BCE. We train once on the
training split (Adam, learning rate $10^{-4}$, weight decay $10^{-5}$, batch 32, horizontal flips), select the
epoch with the lowest development loss, and freeze the checkpoint (SHA-256 recorded). Seeds are fixed (training
seed 42). The frame score is the mean of the map and the video score the mean of 20 frame scores (2 and 19 for the two
documented reduced-frame test videos). Training ran on one NVIDIA A100 GPU and stopped early after \QPADEpochsRun{}
epochs; the selected epoch \QPADBestEpoch{} has development loss \QPADBestDevLoss{} (Fig.~\ref{fig:train}). After a
short single-epoch spike, the development loss flattens while the training loss keeps falling, which is the point
at which early stopping is meant to act.
\qpadfig{training_curves}{Training and development loss per epoch; the dashed line marks the selected
checkpoint.}{fig:train}

\subsection{Threshold analysis and operating point}
A video is bona fide when $\text{score}\ge\tau$. On development scores we enumerate all candidate thresholds,
minimise $|\text{APCER}-\text{BPCER}|$, break ties by the lowest ACER and then the higher threshold. The selected
$\tau=\QPADThreshold$ gives development APCER \QPADDevAPCERatThr{} and BPCER \QPADDevBPCERatThr{}.
This EER operating point balances security and usability without using test data.
\qpadfig{dev_threshold_curve}{Development APCER and BPCER versus threshold; the dashed line is the frozen
operating point.}{fig:thr}

\subsection{Results}
Table~\ref{tab:fp32} reports every rate with its numerator and denominator.
Test print APCER is \QPADTestPrintAPCER, replay APCER \QPADTestReplayAPCER, BPCER \QPADTestBPCER, ACER
\QPADTestACER{} (worst-PAI convention: \QPADTestACERmax), EER \QPADTestEER{} and AUROC \QPADTestAUROC.
Development performance is near perfect (APCER \QPADDevAPCERatThr, BPCER \QPADDevBPCERatThr, AUROC
\QPADDevAUROC), whereas test performance is clearly lower because Protocol~II tests on Printer~2 and Display~2,
which the model never saw. The drop is almost entirely in replay attacks: print APCER stays below 1\%, while
replay APCER rises from 0/270 on development to \QPADTestReplayAPCER{} on test.
George and Marcel report APCER 11.39\%, BPCER 0.56\% and ACER 5.97\% on Protocol~II, with the same worst-PAI APCER
and development-EER threshold~\cite{george2019deep}. Our reproduction reaches a similar ACER (\QPADTestACERmax) with
a lower worst-PAI APCER (\QPADTestAPCERmax) but a higher BPCER (\QPADTestBPCER). Unlike the original paper, which
found print attacks harder than video attacks in most protocols, our errors are dominated by replay. Our recipe
differs from theirs in ways that could explain this: we align faces with the supplied eye coordinates instead of
MTCNN detection and SDM landmarks, we do not undersample the majority class, we use horizontal flips without colour
jitter, and we stop early on development loss instead of training for 50 epochs. We therefore treat the published
result as a sanity reference, not a target.
\begin{table}[!tb]\centering\caption{FP32 DeepPixBiS at the frozen development threshold.}\label{tab:fp32}
\resizebox{\columnwidth}{!}{\qpadtable{fp32_results}}\end{table}

\subsection{Error analysis}
Table~\ref{tab:errors} breaks test errors down by phone and session for each attack type.
Replay attacks account for 25 of the 27 accepted attacks. These errors are concentrated on two phones,
phone~2 (13/60) and phone~4 (8/60), while phones~3 and~6 have none, and they decrease from session~1 (14/120) to
session~3 (3/120). Bona fide rejections are spread evenly (1--3 per phone, 2--4 per session), so no single device
causes them. The test score distribution (Fig.~\ref{fig:testhist}) shows the cause: most attack scores lie near 0 and
most bona fide scores near 1, but a long tail of replay scores extends past the frozen threshold. Replay on unseen
displays is therefore the weakest point of the FP32 baseline, and the attack type for which an INT8-induced change
in APCER has the most room to appear.
\qpadfig{test_score_hist}{Test video scores by class with the frozen development threshold (dashed); attacks to the
right of the threshold are accepted.}{fig:testhist}
\begin{table*}[!t]\centering\caption{Test errors by phone (left) and by session (right) (errors/total).}
\label{tab:errors}\qpadtable{errors_by_phone}\hspace{2.5em}\qpadtable{errors_by_session}\end{table*}

\subsection{Next steps}
\scope{Export the frozen model to ONNX and TensorRT FP32 (score tolerance $10^{-5}$ and $10^{-4}$), audit
calibration scales, and build a six-engine INT8 pilot (natural, no-print, no-replay; two nested seeds) evaluated
with both the reused FP32 threshold and engine-specific development thresholds, before deciding on the full
35-engine study.}

\balance
\bibliographystyle{IEEEtran}
\bibliography{references}
\end{document}
